\section{Results}
\label{sec:results}

Every number in this section is from the \nProjects-project benchmark.
The agents and the retrieval layer are scored on the full \nEval-instance
evaluation split; the two reimplemented published agents, which cost
more to run, are scored on a stratified \nSample-instance sample of that
split (\cref{sec:setup:populations}). All systems in one table are
scored on identical instances. Intervals are 95\% paired bootstrap
intervals over instances; ``significant'' means the interval excludes
zero. Every agent number comes from one run with one seed; what that
costs in precision is measured in \cref{sec:results:negatives}.

\begin{table*}[t]
\caption{File, module and region localization on the full evaluation
split (\nEval{} instances, \nEvalHard{} of them hard, meaning the report
names no file the fix changed). The three systems share the same
evidence pass; the two agents also share the model. Region metrics are
on the \nRegion{} instances whose fix diff yields a gold function.
Cost is per issue: model calls / prompt tokens.}
\label{tab:headline}
\centering
\footnotesize
\setlength{\tabcolsep}{4pt}
\begin{tabular}{@{}lccc cc c c c@{}}
\toprule
 & \multicolumn{3}{c}{file, all} & \multicolumn{2}{c}{file, hard} & module & region & cost \\
\cmidrule(lr){2-4}\cmidrule(lr){5-6}
 & Acc@1 & Acc@5 & MRR & Acc@1 & Acc@5 & Acc@1 & Acc@1 / Acc@5 & calls / tokens \\
\midrule
\tablebody{headline}
\bottomrule
\end{tabular}
\end{table*}

\subsection{What the agent reaches}
\label{sec:results:headline}

\Cref{tab:headline} gives the headline. On \nEval{} real incidents
\sys{} ranks a file the fix changed first for \accOne{} of them and in
the top five for \accFive{} (MRR \mrr). The module it names is correct
for \modOne{}. At the region level, where a hit means the top-ranked
function is one the fix changed, it reaches \regOne{} Acc@1 and
\regFive{} Acc@5 on the \nRegion{} instances that have a gold
function.

The hard stratum is where the challenge puts the difficulty, and it is
half the benchmark. On the \nEvalHard{} reports that name no faulty
file, \sys{} reaches \hardOne{} Acc@1 and \hardFive{} Acc@5; on the easy
half it reaches 83.3\% and 97.4\%. The retrieval layer, which has no
model, reaches \retrOne{} and \retrFive{} overall and 20.3\% and
\retrHardFive{} on the hard stratum. The agent's margin over it is
$+27.4$ points of Acc@1 overall and $+32.5$ on the hard stratum: reading
the code closes most of the gap that the report leaves open.

By report type the picture is flat. Reports with a stack trace score
66.5\% Acc@1, reports with log lines but no trace 66.6\%, and reports
with prose only 70.8\%. A trace does not make the file easier to find at
rank one, because the file it names is usually the one that detected the
failure (\cref{sec:intro}). Across the \nProjects{} projects Acc@1 ranges
from 54.3\% (Infinispan) to 79.6\% (Hadoop). The projects rank the same
way for the retrieval layer (Spearman $\rho = 0.87$ between the two
systems' per-project Acc@1), so project difficulty is largely a property
of the reports rather than of the agent; the per-project table is in the
artifact.

\subsection{What the structure changes}
\label{sec:results:structure}

The unconstrained loop of \cref{sec:setup:systems} ran on the same split
with the same model and the same evidence pass. Paired on the 26{,}537
instances both completed, the staged workflow changes the file-level
result very little: Acc@1 $+0.7$ points $[+0.2, +1.2]$, Acc@5
$-0.9$ $[-1.2, -0.6]$, MRR unchanged. Both differences are significant
and both are small. The loop finds files as well as the workflow does.

What the structure changes is everything below the file. Region Acc@1
rises from \loopRegOne{} to \regOne{} ($+14.1$ points $[+13.5, +14.7]$)
and region Acc@5 from \loopRegFive{} to \regFive{} ($+19.2$
$[+18.6, +19.9]$). The loop names 3.2 regions per answer, one function
per file; the workflow judges every examined file with its functions on
screen and names up to three of them. The explanation changes in kind as well: in
the loop's chains no step cites a log event, because nothing required
it; in the workflow 57.4\% of chain steps cite one of the report's
selected log events, and 85.7\% of those citations exist in the report
word for word. The price is \calls{} model calls and \tokens{} prompt
tokens per issue against the loop's 13.6 and 99k, and 286 against 199
seconds of wall time at the same concurrency.

By report type the file-level change is nil on reports with a trace
($-0.3$ $[-1.3, +0.7]$) and small but positive on reports with logs only
($+1.4$ $[+0.3, +2.4]$) and prose only ($+1.1$ $[+0.3, +1.8]$). On the
hard stratum the workflow is $0.5$ points below the loop at Acc@1
$[-1.3, +0.3]$, not significant, and $1.1$ below at Acc@5
$[-1.7, -0.5]$, while its region Acc@1 is $+13.4$ $[+12.5, +14.2]$ and
region Acc@5 $+20.9$ $[+19.9, +21.9]$ above.

\subsection{Against reimplemented published agents}
\label{sec:results:agents}

\begin{table}[t]
\caption{On the \nSample-instance stratified sample (\nSampleHard{}
hard), all systems completed, all on the same 27B model. Hard columns
are Acc@1 and Acc@5 on the hard stratum; cost is model calls and prompt
tokens per issue. LocAgent-style and CoSIL-style are our
reimplementations of the two published agents on our tree-sitter layer;
``model alone'' is the two-stage Agentless procedure over the file
tree; the re-ranker reorders the retrieval layer's top twenty
(\cref{sec:setup:systems}).}
\label{tab:agents}
\centering
\footnotesize
\setlength{\tabcolsep}{3pt}
\begin{tabular}{@{}lcccccc@{}}
\toprule
system & Acc@1 & Acc@5 & hard@1 & hard@5 & calls & tokens \\
\midrule
\tablebody{agents}
\bottomrule
\end{tabular}
\end{table}

\Cref{tab:agents} compares \sys{} with two agents from the literature,
reimplemented because neither released system runs on Java, Scala, C or
C++ (\cref{sec:setup:systems}), and with the same model used without
tools, alone and as a re-ranker of the retrieval layer's candidates.
\sys{} leads every system on every column. Against LocAgent-style, the
stronger published agent, the paired difference is $+0.072$ Acc@1
$[+0.045, +0.100]$ and $+0.136$ Acc@5 $[+0.110, +0.162]$; on the hard
stratum $+0.118$ $[+0.074, +0.164]$ and $+0.197$ $[+0.153, +0.240]$.
Against CoSIL-style it is $+0.199$ and $+0.257$ overall and $+0.350$
and $+0.443$ on the hard stratum. Against the model alone it is
$+0.144$ Acc@1 $[+0.113, +0.176]$ and $+0.141$ Acc@5
$[+0.116, +0.167]$; against the re-ranker, the strongest baseline at
rank one, $+0.063$ $[+0.038, +0.088]$ and $+0.061$ $[+0.042, +0.082]$.
Every difference is significant. Because all five systems share the
model, what separates them is the architecture.

Cost reads the other way. The re-ranker reaches 0.608 Acc@1 for one
model call and 2.8k prompt tokens per issue, a fortieth of what \sys{}
spends, and it is the right choice where that budget is the
constraint. LocAgent-style spends the same tokens as \sys{} (106k
against 108k) for seven points less at rank one; CoSIL-style spends a
quarter of the tokens and loses most of the hard stratum. The
re-ranker's limit is its candidate set: it can only reorder the twenty
files retrieval found, and on the hard stratum it falls to 0.393
against \sys{}'s 0.522.

Two observations about the ports. LocAgent-style called its graph
traversal tool 0.09 times per run against 14.5 keyword searches, so the
multi-hop traversal that distinguishes LocAgent in its own paper goes
almost unused under this model. And both arms returned no parseable
location for 8\% to 10\% of instances (LocAgent-style 8.2\%, CoSIL-style
10.1\%). CoSIL's reflective alignment step exists to repair exactly that
failure, so we expected CoSIL-style to show the lower rate; it does not,
and we record that as a property of our port we cannot explain rather
than as a finding about CoSIL. \sys{} returned an answer on every
instance, because its submission is a schema the run cannot end
without.

\textbf{The model does not remember these fixes.} These are public
projects, and the model may simply have memorised their histories. We
asked it directly on the same \nSample{} instances: given only an issue
key such as KAFKA-6624 and no report text, which files did that fix
touch? It claimed to recall nothing and scored 0.0\% Acc@5. Given the
issue title as well, it claimed recall for 2.1\% of instances and scored
1.6\%. What the model-based systems contribute is inference from the
report and the code, not recitation.

\subsection{The retrieval layer, and what memory is worth}
\label{sec:results:retrieval}

\begin{table}[t]
\caption{The retrieval layer's signals and their fusion, pooled over the
\nProjects{} projects on the evaluation split (\nEval{} instances;
\nEvalHard{} hard). ``Conditioned'' is the regime-conditioned weighted
fusion of \cref{sec:approach:fusion}.}
\label{tab:retrieval}
\centering
\scriptsize
\setlength{\tabcolsep}{2.5pt}
\begin{tabular}{@{}lccc cc c@{}}
\toprule
 & \multicolumn{3}{c}{all} & \multicolumn{2}{c}{hard} & easy \\
\cmidrule(lr){2-4}\cmidrule(lr){5-6}
 & Acc@1 & Acc@5 & MRR & Acc@1 & Acc@5 & Acc@5 \\
\midrule
\tablebody{retrieval}
\bottomrule
\end{tabular}
\end{table}

\sys{} uses the retrieval layer as a tool, so \cref{tab:retrieval}
reports how that layer behaves on its own across the whole benchmark.
The signals divide cleanly by stratum. Trace mentions (S1) are the best
single signal on easy reports and score exactly zero on hard ones, by
definition. Content retrieval (S3) is the best standalone signal
overall. Incident memory (S4) is the one signal whose accuracy barely
moves between the strata (43.8\% against 40.2\% Acc@5), because a past
fix can point at a file the present report never names, and on the hard
stratum it is the second-best standalone signal.

That is where memory earns its place. Adding it to unweighted fusion
lifts hard-stratum Acc@5 from 26.4\% to 41.8\%; adding it to the
conditioned fusion, which already down-weights the mention signals in
the no-mention regime, lifts it from 46.5\% to 48.7\%. The increment is
small once the weights are tuned per regime and large when they are
not, which is the practical reading: memory is cheap insurance against
a report that names nothing.

The memory signal is evaluated behind the firewall of
\cref{sec:approach:firewall}. Pooled over the projects, the
future-reports rule alone removes \firewallPct{} of the candidate votes
the index proposes, and the numbers above are what survives that cut.

\subsection{Regions and modules}
\label{sec:results:regions}

\begin{table}[t]
\caption{Function-level results on the evaluation instances whose fix
diff yields at least one gold function. ``Oracle'' is the share whose
gold function is anywhere in the candidate set; it is the ceiling of
the BM25 stage and is not defined for the agents, which choose their
own regions.}
\label{tab:functions}
\centering
\scriptsize
\setlength{\tabcolsep}{2.5pt}
\begin{tabular}{@{}lcccccc@{}}
\toprule
 & eligible & Acc@1 & Acc@5 & Acc@10 & MRR & oracle \\
\midrule
\tablebody{functions}
\bottomrule
\end{tabular}
\end{table}

The challenge asks for code regions, not files. Gold functions are
derived mechanically: tree-sitter lists every function span in the
snapshot, and a function is gold when its span overlaps a hunk of the
fix diff; 93.9\% of instances yield at least one. \Cref{tab:functions}
compares three ways of getting there. A BM25 stage that ranks the
function bodies in the retrieval layer's top ten files reaches 21.9\%
Acc@1 with a 77.0\% ceiling, so its problem is ranking inside large
files, not finding them. The unconstrained loop reaches 39.8\%, and \sys{}
53.8\%, because its \textsc{examine} step reads each candidate file
whole, or in whole functions, and names the functions that hold the
fault while the code is on screen. On the hard stratum \sys{} reaches
39.2\% region Acc@1 and 55.2\% Acc@5.

At the module level, where modules are the leading path components of a
file at a fixed per-project depth, \sys{} names the correct module for
\modOne{} of instances (85.5\% hard, 94.5\% easy), and the module of its
top-ranked file is correct for 90.9\%. The retrieval layer's top file
gives \moduleRetr{}. Module spaces are small and defects concentrate, so
these numbers need the sanity baselines of \cref{tab:modules}: ranking
the modules that exist in the snapshot by their defect frequency on the
development split already reaches \moduleFreq{} Acc@1, and the agent's
margin over that prior, 41 points, is what counts as localization.

\begin{table}[t]
\caption{Module level on the evaluation split, with two sanity
baselines: the expected accuracy of a random module ranking
(analytical, per snapshot) and ranking the snapshot's modules by their
defect frequency on the development split.}
\label{tab:modules}
\centering
\footnotesize
\begin{tabular}{@{}lcc@{}}
\toprule
 & Acc@1 & Acc@5 \\
\midrule
\tablebody{modules}
\bottomrule
\end{tabular}
\end{table}

\subsection{Cost, and a dense alternative}
\label{sec:results:cost}

\sys{} spends \calls{} model calls and \tokens{} prompt tokens per
issue and 286 seconds of wall time at 32 concurrent issues on one H100,
which came to about six issues per minute per GPU over the run. The
retrieval layer answers from one CPU core with no model, in a median of
under 0.2 seconds per issue on every project. That
is a gap of four orders of magnitude for $+27.4$ points of Acc@1, and
which side of it is right depends on whether a team can host a 27B
model. A team that cannot, because source and production logs may not
leave the building, still has a pipeline with no GPU, no training and no
network that puts a changed file in the top five for seven issues in
ten.

The challenge asks for joint log and code semantics, so we also measured
what a small dense retriever buys inside the retrieval layer. On
200-instance samples per project, replacing BM25 with bge-small
embeddings in the content slot loses on every one of the \nProjects{}
projects (58.4\% to 44.9\% mean Acc@5): matching a log line to code is
mostly identifier matching, and the embedder reads at most 1,400 words
of a file where BM25 sees all of it. In the incident-memory slot, where
the task is report-to-report similarity, dense retrieval wins on 22 of
27 projects (37.6\% to 40.4\%) and lifts the full system from 67.7\% to
69.1\% with weights tuned for the lexical signal. Semantics pays in the
memory slot first; the code slot needs the model to read the code, which
is what the agent does.

\subsection{What goes against us}
\label{sec:results:negatives}

Five things. First, the staged workflow is below the unconstrained loop at
Acc@5 by 0.9 points, significantly, and gains only 0.7 at Acc@1: at the
file level the structure buys nothing, and a reader who wants files only
should take the loop. Second, on reports with a stack trace the
workflow is not better at all. Third, four of \nEval{} instances
(0.015\%) ended without an answer in two runs each, because the final
step's structured output failed to parse twice; they are scored as
misses and not retried. Fourth, the port check for CoSIL-style fails,
as stated above. Fifth, the agent is not deterministic and we report one
seed. On 20 instances run twice with the same seed on the same machine,
the top-ranked file agreed on 18 and the top ten on 12; across two
machines with identical software and model checksums, the top file
agreed on 14 to 15 and the module on all 20, and Acc@1 moved between
0.65 and 0.70. A difference below about two points of Acc@1 between two
agent runs should be read as noise, which is why the file-level
differences in \cref{sec:results:structure} are reported as small even
where they are significant.
